\chapter{Introduction}
\label{chap:introduction}

The rapid convergence of Large Language Models (LLMs) and Multi-Agent Systems (MAS) has initiated a paradigm shift in the simulation of complex socio-political and economic dynamics. What began as experimental role-playing within unstructured textual environments has rapidly evolved into the deployment of autonomous agents capable of long-term strategic reasoning, negotiation, and resource management. This thesis explores the frontier of geopolitical simulation by introducing a novel framework designed to rigorously evaluate the behavior of LLM-driven state actors under conditions of scarcity, imperfect information, and ideological friction.

\section{Context}
\label{sec:intro_context}
The integration of artificial intelligence into statecraft and military strategy is no longer a theoretical exercise isolated to academic sandboxes. As demonstrated by recent real-world deployments, ranging from the Albanian government's introduction of ``Diella,'' the first AI-generated cabinet minister governing public procurement alongside a network of parliamentary sub-agents, to the highly publicized institutional disputes between the U.S. Pentagon and AI laboratories over the deployment of models like Claude on classified military networks, LLMs are actively transitioning from analytical tools to operational strategic agents.
As these models become embedded in the command-and-control infrastructure of global superpowers, simulating and understanding their emergent behaviors, deterrence strategies, and propensities for conflict escalation becomes an urgent scientific and security imperative. 

\section{Motivations and Objectives}
\label{sec:intro_motivations}
Despite the accelerated adoption of LLMs in wargaming and policy simulation, the current State of the Art exhibits structural limitations that hinder verifiable scientific inquiry. Existing monolithic frameworks often operate as opaque "black boxes" (e.g., Cicero) or produce unstructured "stream of consciousness" traces (e.g., WarAgent) that resist quantitative causal analysis. Furthermore, these environments frequently lack a deterministic "physics engine," allowing agents to hallucinate actions unconstrained by material or topological realities. Crucially, when tested on real-world maps, modern LLMs suffer from severe "historical data leakage," reverting to pre-trained historical memorization rather than demonstrating autonomous strategic reasoning.

The primary objective of this thesis is to address these empirical gaps by designing and validating \textbf{GeoMAS} (Geopolitical Multi-Agent Simulation). GeoMAS is a neuro-symbolic framework that establishes a \textit{Duality of Determinism}: it physically grounds the stochastic, hierarchical cognitive reasoning of LLM agents (representing national cabinets) within a mathematically rigid, procedurally generated environment (Voronoi geography). By decoupling strategic agency from historical bias and enforcing strict resource and topological constraints through a deterministic ``Rules Oracle,'' the framework aims to provide a verifiable sandbox for studying emergent geopolitical equilibria, the weaponization of diplomatic rhetoric, and the macroscopic impact of exogenous shocks.

\section{Main Results and Contributions}
\label{sec:intro_results}
Through extensive experimental validation, this thesis demonstrates that the GeoMAS framework successfully models complex, interpretable macro-structures without degrading into chaotic noise. The primary contributions and findings of this research include:

\begin{itemize}
    \item \textbf{Coalition-Based Deterrence and Bipolarity:} The baseline simulations demonstrate that the system does not converge to a classical, symmetric multipolar balance of power, nor into chaotic unconstrained warfare. Instead, it systematically self-organizes into a structurally bipolar configuration: a heavily integrated, NATO-style \textbf{Mutual Defense Coalition} of four democracies coordinating for collective security, counterbalanced by an isolated, heavily militarized authoritarian state that accumulates latent power without initiating conflict. Active warfare is relegated entirely to peripheral, non-aligned states.
    \item \textbf{Strategic Deception and Moral Washing:} By isolating the agents' \textit{private} strategic intents from their \textit{public} diplomatic declarations, the framework successfully operationalizes and quantifies LLM deception. The analysis reveals that "Global Strategy" strongly dominates "Government Type" in predicting deceptive behavior. Notably, when aggressive expansionist doctrines conflict with cooperative ideologies (e.g., Democratic regimes), the agents do not moderate their aggression; rather, they resolve the cognitive dissonance through sophisticated \textit{Moral Washing}, explicitly weaponizing defensive and democratic rhetoric to mask unprovoked empirical conquest.
\end{itemize}

% NOTE: The following paragraph is explicitly designed to be easily expanded or updated 
% once the quantitative and qualitative analyses for RQ3 (Scenarios) and RQ4 (Counterfactuals/XAI) are finalized.

\section{Structure of the Thesis}
\label{sec:intro_structure}
The remainder of this thesis is organized as follows:
\begin{itemize}
    \item \textbf{Chapter 2 (State of the Art in Geopolitical Simulation):} Reviews the evolution of LLM-based cognitive architectures, negotiation protocols, safety risks in wargaming, and identifies the core limitations bridging simulation to real-world deployment.
    \item \textbf{Chapter 3 (The GeoMAS Framework):} Details the neuro-symbolic architecture of the simulation, including the hierarchical cabinet model, the procedural Voronoi environment, the deterministic Rules Oracle, and the memory management systems.
    \item \textbf{Chapter 4 (Experimental Methodology and Validation):} Outlines the structured experimental design, defining the baseline parameters, the tracked telemetry, the scenario injection protocols, and the specific Research Questions guiding the validation.
    \item \textbf{Chapter 5 (Results):} Presents the core empirical findings derived from the baseline simulations, analyzing emergent sociopolitical equilibria (RQ1) and microdynamic strategic deception (RQ2).
    \item \textbf{Chapter 6 (Threats to Validity):} Critically assesses the methodological limits of the research, discussing the boundaries of LLM determinism, the heuristic biases of context engineering, and the generalizability of results derived from a single model and topology.
    \item \textbf{Chapter 7 (Conclusion):} Summarizes the primary contributions of the thesis, reflects on the broader implications for the development of verifiable AI governance tools, and proposes avenues for future research.
\end{itemize}